\section{Conclusion}
\label{sec:conclusion}

We have presented a formal framework for reasoning about the long-term
viability of AI strategies in competitive, multi-agent environments. By
leveraging concepts from evolutionary game theory, we identified
a class of strategies---potential-maximizing strategies---that are
evolutionarily stable and uninvadable.

Our central result is that any universe possessing a horizon admits a
potential-maximizing strategy, and that under fairness (or the weaker economies
of scale assumption), such a strategy maintains or grows its share of potential
over time. Crucially, a potential-maximizing strategy never actually maximizes
its reward directly; instead, it perpetually invests in maximizing its
\emph{potential} to do so. This behavior is strikingly consistent with the
instrumental convergence thesis~\citep{bostrom2012} and the formalization of
power-seeking behavior~\citep{powerseeking}: regardless of the terminal goal, a
robust strategy converges on resource acquisition, self-preservation, and
capability growth.

For AI safety, these findings carry a sobering implication. An aligned strategy
that is not potential-maximizing risks being outcompeted and displaced by one
that is. Alignment techniques must therefore be evaluated not only for their
correctness but for their competitive fitness---their ability to persist in a
landscape where rival strategies are under constant selective pressure to
improve. Our framework offers a formal lens through which to conduct such
evaluations.
